(1) $\Cl(H_0)$ has dimension $4$; the elements $1, v, w_0$ are independent (grades $0$ and $1$), and $v w_0 = v \cdot w_0 + v \wedge w_0 = \tfrac12 + v \wedge w_0$ has a nonzero grade-$2$ component, so the four are a basis. Every vector of $U'$ is orthogonal to $H_0$, hence anticommutes with $v$ and $w_0$; consequently any product of generators of $\Cl(\mathbb{V}')$ can be reordered, changing only signs, into (an element of $\Cl(H_0)$)$\,\times\,$(an element of $\Cl(U')$). The products of the four basis elements of $\Cl(H_0)$ with the $2^{\dim U'}$ basis blades of $\Cl(U')$ therefore span $\Cl(\mathbb{V}')$, and since their number $4 \cdot 2^{\dim U'} = 2^{\dim \mathbb{V}'}$ equals $\dim \Cl(\mathbb{V}')$ (valid whether or not $Q'|_{U'}$ is degenerate, by the count of Section \ref{sec:blades} \cite{Chevalley1997}), they form a basis. This is the graded tensor decomposition $\Cl(A \perp B) \cong \Cl(A)\,\hat{\otimes}\,\Cl(B)$ \cite{Lam2005} in coordinates.

(2) For $x \in \Cl(U')$, the element $vx$ \emph{is} an expansion of the form in (1), with $\beta = x$ and $\alpha = \gamma = \delta = 0$. By uniqueness, $vx = 0$ forces $x = 0$.
